\section{Background Literature}\label{backg}

This section reviews the literature that informed the project, in four strands: the state of automated programming assessment, the field's persistent weakness in feedback, the design principles used to evaluate and improve the interfaces, and the standards and practices that shaped the engineering. Each strand ends where the project begins, and the review closes by positioning the project against them.

\subsection{Automated Assessment of Programming}

Assessing programming by hand does not scale, and it barely works even when it is affordable: Lobb and Harlow open the CodeRunner paper by observing that hand-written code on paper is nearly impossible to grade meaningfully, because static reading cannot reveal whether a student could make their program actually run~\cite{CodeRunner2016}. The field's answer, developed over decades, is to execute submissions against instructor-defined test cases and grade the outcome. Ihantola et al.'s systematic review of assessment systems shows both the maturity and the fragmentation of this space: dozens of tools exist, most built around the same test-based core, differing in how teachers define tests, how resubmission is handled, and how securely untrusted code is run~\cite{reviewrecentsystems2010}. Two of that review's conclusions matter here. First, test-based assessment is the settled approach; CodeRunner is an orthodox member of the family. Second, the authors explicitly criticise the field's habit of building yet another new system, and encourage opening up and extending the systems that already exist. That recommendation is a direct argument for the shape of this project: improving a widely deployed tool rather than producing a new one.

Pieterse's analysis of automated assessment adds the operational perspective: the success of such systems depends on factors beyond grading correctness, including how students experience failure and how much confidence teachers have in the tool, and the requirements placed on these systems grow as education leans harder on them~\cite{pieterse2013automated}. A tool can therefore be complete in capability yet still fall short of what its users need, which is precisely the position this project found CodeRunner in.

\subsection{The Feedback Problem}

The most consistent criticism in the assessment literature concerns feedback. Keuning, Jeuring and Heeren's systematic review of automated feedback generation finds that most tools report whether a submission is correct but offer little help with why it failed or how to move forward~\cite{keuningfeedback}. This is the academic form of the complaint that surfaced independently in this project's survey: hidden test cases fail without telling the student what went wrong, leaving them to guess at the cost of further penalty. The literature thus both anticipates the pain point and elevates it from an anecdote to a known, field-wide weakness, which is why hints for hidden test cases (FR5) is treated as the highest-value item of future work rather than a passing suggestion.

A newer strand asks whether large language models can close the feedback gap. Phung et al.\ benchmark ChatGPT and GPT-4 against human tutors across programming education tasks, including grading, and find that even the strongest models approach but do not match human tutors~\cite{llmgrading}. For grading specifically, where consistency is a fairness requirement, approaching is not enough; this evidence underpins the decision, discussed in Section~\ref{process}, to consider and exclude LLM-based grading from the project's scope.

\subsection{Usability and Design Principles}

The project's central claim, that CodeRunner's remaining problems are experience problems, rests on established design literature. Norman's \emph{The Design of Everyday Things} argues that what is habitually called user error is usually design error: systems that punish slips, hide their state, or demand documentation to operate are failing their users, not the reverse~\cite{designofeveryday}. Norman also supplies the vocabulary of affordances, controls whose appearance invites their use, which this project applies directly in the design of its student-facing controls. Krug's \emph{Don't Make Me Think} distils the same stance for the web: an interface should be self-evident, and every moment a user spends puzzling over a control is a cost paid for a design decision~\cite{dontmakemethink}. These two works served as the lenses for the heuristic review of CodeRunner's pages reported in Section~\ref{process}.

Written guidance is part of the interface as well. Google's developer documentation style guide asks writers to organise content around the reader's task, keep pages short and scannable, and explain no more than the task requires~\cite{googledevstyle}. It was used as the yardstick for CodeRunner's documentation, and the requirement it produced (NFR3) is essentially the guide's advice turned into an acceptance criterion.

\subsection{Standards and Engineering Practice}

Accessibility work in this project is anchored to two W3C standards: the Web Content Accessibility Guidelines define what an accessible page must achieve~\cite{w3c_wcag22_2023}, and WAI-ARIA defines the vocabulary of attributes used to achieve it in rich interfaces, including the labelling and live-region techniques applied here~\cite{w3c_wai_aria12_2023}. The Lighthouse audit used to measure progress derives its accessibility checks from these guidelines, which makes the requirement NFR2 measurable rather than aspirational.

The engineering practices around the product likewise follow the literature. Docker's containerisation, packaging applications with their dependencies into isolated and reproducible units~\cite{dockerjournal}, is what makes the development environment requirements (DE1 and DE2) satisfiable with one command. Behaviour-driven development, as introduced by North, frames acceptance tests as executable descriptions of behaviour~\cite{northbdd}; CodeRunner's upstream Behat suite is exactly such a description, and running it unmodified is this project's regression evidence. The survey instrument was designed with reference to Kitchenham and Pfleeger's principles of survey research in software engineering~\cite{designingsurvey}, and the validation of requirements against implementations follows the traceability practice analysed by Gotel and Finkelstein~\cite{goteltraceability}.

\subsection{Positioning of this Project}

Read together, the strands position the project precisely. The assessment literature says the grading problem is solved and advises improving existing systems over building new ones; the feedback literature identifies experience, not capability, as the open weakness; the design literature supplies principles and methods for fixing experience problems and insists such problems belong to the software; and the standards make the improvements measurable. This project sits at that intersection: an evidence-driven user experience improvement of a mature, widely used automated assessment tool, evaluated with its real users.
